\chapter{Reaction Mechanisms: Curly Arrows}\label{ch:g12:curly-arrows}

A \omterm{def:g8:balanced-equations:equation}{balanced equation} is like the first and last photographs of a film: it
shows the \omterm{def:g7:chemical-reactions:reactant}{reactants} before and the products after, and nothing in
between. Yet \omterm{def:g7:atoms-and-molecules:molecule}{molecules} do not simply swap \omterm{def:g7:atoms-and-molecules:atom}{atoms}. Bonds break and form one
at a time, as pairs of \omterm{def:g9:inside-the-atom:nucleus}{electrons} move from one place to another, through
short-lived species that never appear in the equation. Chemists draw
these movements with curved arrows. With a few rules, these arrows
explain why a reaction happens where it does, and predict reactions
never seen before.

\begin{recall}
\omterm{def:g11:polarity-and-cohesion:electronegativity}{Electronegativity}, \omterm{def:g11:polarity-and-cohesion:polar-bond}{polar bonds} and \omterm{def:g11:polarity-and-cohesion:polar-bond}{partial charges}
(\cref{ch:g11:polarity-and-cohesion}). \omterm{def:g10:lewis-and-shape:lewis-structure}{Lewis structures}: \omterm{def:g10:lewis-and-shape:lone-pair}{bonding pairs} and
\omterm{def:g10:lewis-and-shape:lone-pair}{lone pairs} (\cref{ch:g10:lewis-and-shape}). \omterm{def:g11:functional-groups:functional-group}{Functional groups} and
families (\cref{ch:g11:functional-groups}).
\end{recall}

\section{Electron-rich and electron-poor sites}

\begin{definition}[Electron-donor and electron-acceptor sites]\label{def:g12:curly-arrows:site}
An \emph{electron-donor site}\index{electron-donor site} of a \omterm{def:g7:atoms-and-molecules:molecule}{molecule} or
an \omterm{def:g9:ions:ion}{ion} is a place rich in \omterm{def:g9:inside-the-atom:nucleus}{electrons}, able to give a pair to form a new
bond: a \omterm{def:g10:lewis-and-shape:lone-pair}{lone pair}, a \omterm{def:g10:lewis-and-shape:multiple-bond}{double bond}, an \omterm{def:g7:atoms-and-molecules:atom}{atom} carrying a negative charge or a
\omterm{def:g11:polarity-and-cohesion:polar-bond}{partial charge} $\delta^-$. An \emph{electron-acceptor
site}\index{electron-acceptor site} is a place poor in \omterm{def:g9:inside-the-atom:nucleus}{electrons}, able to
receive a pair: an \omterm{def:g7:atoms-and-molecules:atom}{atom} carrying a positive charge or a \omterm{def:g11:polarity-and-cohesion:polar-bond}{partial charge}
$\delta^+$.
\end{definition}

\begin{example}[Sites in two molecules]\label{ex:g12:curly-arrows:sites}
In bromoethane, \ce{CH3-CH2-Br}, bromine is more electronegative than
carbon: the carbon bonded to it carries $\delta^+$ (acceptor site), the
bromine $\delta^-$ and three \omterm{def:g10:lewis-and-shape:lone-pair}{lone pairs}. In propanone, the oxygen of the
\ce{C=O} group carries $\delta^-$ and two \omterm{def:g10:lewis-and-shape:lone-pair}{lone pairs} (donor sites), its
carbon $\delta^+$ (acceptor site). The hydroxide \omterm{def:g9:ions:ion}{ion} \ce{OH-}, with its
negative charge and three \omterm{def:g10:lewis-and-shape:lone-pair}{lone pairs}, is a strong donor.
\end{example}

\begin{omfigure}
\setchemfig{atom sep=2.2em}
\begin{tabular}{cc}
\chemfig{H_3C-C(-[2]H)(-[6]H)-\charge{90=\:,0=\:,270=\:}{Br}} &
\chemfig{H_3C-C(-[6]CH_3)=[2]\charge{0=\:,180=\:}{O}} \\[2pt]
\small $\delta^+$ on the carbon bonded to \ce{Br}, $\delta^-$ on \ce{Br} &
\small $\delta^+$ on the carbon of \ce{C=O}, $\delta^-$ on \ce{O} \\
\small bromoethane & \small propanone
\end{tabular}

{\small Electron-poor carbon \omterm{def:g7:atoms-and-molecules:atom}{atoms} ($\delta^+$, acceptor sites) next to
electron-rich, electronegative \omterm{def:g7:atoms-and-molecules:atom}{atoms} ($\delta^-$, with \omterm{def:g10:lewis-and-shape:lone-pair}{lone pairs}: donor
sites).}
\end{omfigure}

\section{Curly arrows}

\begin{definition}[Curly arrow]\label{def:g12:curly-arrows:curly-arrow}
A \emph{curly arrow}\index{curly arrow} shows the movement of a pair of
\omterm{def:g9:inside-the-atom:nucleus}{electrons} during a step of a reaction. It starts from the pair that
moves, a \omterm{def:g10:lewis-and-shape:lone-pair}{lone pair} or a bond, and ends where the pair goes: on an \omterm{def:g7:atoms-and-molecules:atom}{atom},
where it forms a new bond, or on an \omterm{def:g7:atoms-and-molecules:atom}{atom} that takes the pair of a bond
that breaks.
\end{definition}

\begin{method}[Drawing curly arrows]\label{met:g12:curly-arrows:drawing}
\begin{enumerate}
  \item Find the donor site and the acceptor site.
  \item Draw an arrow from the donor pair (a \omterm{def:g10:lewis-and-shape:lone-pair}{lone pair} or a bond) to the
        acceptor \omterm{def:g7:atoms-and-molecules:atom}{atom}: a new bond forms there.
  \item If the acceptor \omterm{def:g7:atoms-and-molecules:atom}{atom} would then have too many \omterm{def:g9:inside-the-atom:nucleus}{electrons} (more
        than an octet, or more than two for hydrogen), draw a second arrow
        breaking one of its bonds, the pair going to the more
        electronegative \omterm{def:g7:atoms-and-molecules:atom}{atom}.
  \item Check the charges after the step: the \omterm{def:g7:atoms-and-molecules:atom}{atom} that gave a \omterm{def:g10:lewis-and-shape:lone-pair}{lone pair}
        becomes more positive by one, the \omterm{def:g7:atoms-and-molecules:atom}{atom} that received a pair as a
        \omterm{def:g10:lewis-and-shape:lone-pair}{lone pair} more negative by one; the total charge does not change.
\end{enumerate}
\end{method}

\begin{proposition}[Charge is conserved in each step]\label{prop:g12:curly-arrows:charge}
The total electric charge of the species is the same before and after
each step of a mechanism.
\end{proposition}

\begin{proof}
A \omterm{def:g12:curly-arrows:curly-arrow}{curly arrow} only moves \omterm{def:g9:inside-the-atom:nucleus}{electrons} that are already there; no \omterm{def:g9:inside-the-atom:nucleus}{electron}
is created or destroyed, and the nuclei do not change.
\end{proof}

\begin{example}[A first arrow]\label{ex:g12:curly-arrows:first}
A hydrogen \omterm{def:g9:ions:ion}{ion} \ce{H+} meets a hydroxide \omterm{def:g9:ions:ion}{ion} \ce{OH-}. A \omterm{def:g10:lewis-and-shape:lone-pair}{lone pair} of the
oxygen forms a bond with the hydrogen \omterm{def:g9:ions:ion}{ion}, which has no \omterm{def:g9:inside-the-atom:nucleus}{electrons}:
\ce{H+ + OH- -> H2O}. One arrow, from the \omterm{def:g10:lewis-and-shape:lone-pair}{lone pair} of the oxygen to the
\ce{H+}; the charges $+1$ and $-1$ become 0 in the water \omterm{def:g7:atoms-and-molecules:molecule}{molecule}.
\end{example}

\begin{omfigure}
\setchemfig{atom sep=2.1em}
\chemfig{@{hp}H^{+}}\qquad\raisebox{0.2ex}{$+$}\qquad
\chemfig{H-@{ox}\charge{90=\:,0=\:,270=\:}{O}^{-}}
\qquad\raisebox{0.2ex}{$\longrightarrow$}\qquad
\chemfig{H-\charge{90=\:,270=\:}{O}-H}
\chemmove{
  \draw[->, thick, omProp, shorten <=4pt, shorten >=3pt]
    ($(ox)+(0,0.25)$) .. controls +(110:9mm) and +(60:9mm) .. (hp);}

{\small The simplest step: a \omterm{def:g10:lewis-and-shape:lone-pair}{lone pair} of the hydroxide \omterm{def:g9:ions:ion}{ion} moves to form
the new \ce{O-H} bond with the hydrogen \omterm{def:g9:ions:ion}{ion}.}
\end{omfigure}

\section{Elementary steps and intermediates}

\begin{definition}[Mechanism, elementary step, intermediate]\label{def:g12:curly-arrows:mechanism}
The \emph{reaction mechanism}\index{reaction mechanism} of a reaction is
the series of steps by which the \omterm{def:g7:chemical-reactions:reactant}{reactants} become the products. Each
\emph{elementary step}\index{elementary step} is a single event in which
a few pairs of \omterm{def:g9:inside-the-atom:nucleus}{electrons} move at once. A species formed in one step and
used up in a later one is a \emph{reaction intermediate}\index{reaction
intermediate}: it does not appear in the overall equation.
\end{definition}


\begin{example}[A carbocation intermediate]\label{ex:g12:curly-arrows:carbocation}
When hydrogen bromide adds to ethene, the reaction takes two steps. In
the first, the \omterm{def:g10:lewis-and-shape:multiple-bond}{double bond} of ethene gives a pair to the hydrogen of
\ce{HBr}, and the \ce{H-Br} pair goes to bromine, which leaves as
\ce{Br-}; the carbon that lost its share of the \omterm{def:g10:lewis-and-shape:multiple-bond}{double bond} is left with
only six \omterm{def:g9:inside-the-atom:nucleus}{electrons} and a positive charge: a carbocation,
\ce{CH3-CH2+}. In the second step, a \omterm{def:g10:lewis-and-shape:lone-pair}{lone pair} of \ce{Br-} forms a bond
with that carbon. The carbocation is an intermediate.
\end{example}

\section{Three categories of reaction}

\begin{definition}[Substitution, addition, elimination]\label{def:g12:curly-arrows:categories}
In a \emph{substitution}\index{substitution}, an \omterm{def:g7:atoms-and-molecules:atom}{atom} or group of a
\omterm{def:g7:atoms-and-molecules:molecule}{molecule} is replaced by another. In an \emph{addition}\index{addition},
two \omterm{def:g7:atoms-and-molecules:molecule}{molecules} join into one, a multiple bond becoming a single one. In
an \emph{elimination}\index{elimination}, a small \omterm{def:g7:atoms-and-molecules:molecule}{molecule} (often water)
is removed from a \omterm{def:g7:atoms-and-molecules:molecule}{molecule}, a single bond becoming a multiple one.
\end{definition}

\begin{omfigure}
\setchemfig{atom sep=2em}
\small
\textbf{\omterm{def:g12:curly-arrows:categories}{Substitution}} (one step): \ce{CH3-CH2-Br + OH- -> CH3-CH2-OH + Br-}

\medskip
\chemfig{H-@{o1}\charge{90=\:,0=\:,270=\:}{O}^{-}}\qquad
\chemfig{H_3C-@{c1}C(-[2]H)(-[6]H)-[@{b1}]@{br1}\charge{90=\:,0=\:,270=\:}{Br}}
\qquad$\longrightarrow$\qquad
\chemfig{H_3C-C(-[2]H)(-[6]H)-O-H}\quad$+$\quad\ce{Br-}
\chemmove{
  \draw[->, thick, omProp, shorten <=4pt, shorten >=3pt]
    ($(o1)+(0.15,0.2)$) .. controls +(60:9mm) and +(120:9mm) .. (c1);
  \draw[->, thick, omDef, shorten <=2pt, shorten >=4pt]
    (b1) .. controls +(-90:7mm) and +(-90:7mm) .. (br1);}

\vspace{6mm}
\textbf{\omterm{def:g12:curly-arrows:categories}{Addition}} (two steps): \ce{CH2=CH2 + HBr -> CH3-CH2-Br}

\vspace{9mm}
\chemfig{H_2C=[@{pi}]CH_2}\quad$+$\quad
\chemfig{@{h2}H-[@{b2}]@{br2}Br}
\qquad$\longrightarrow$\qquad
\chemfig{H_3C-@{cp}\charge{90:2pt=$\scriptstyle +$}{C}H_2}\quad$+$\quad
\chemfig{@{brm}\charge{90=\:,180=\:,270=\:,0=\:}{Br}^{-}}
\qquad$\longrightarrow$\qquad \ce{CH3-CH2-Br}
\chemmove{
  \draw[->, thick, omProp, shorten <=2pt, shorten >=3pt]
    (pi) .. controls +(90:10mm) and +(110:9mm) .. (h2);
  \draw[->, thick, omDef, shorten <=2pt, shorten >=4pt]
    (b2) .. controls +(-90:7mm) and +(-90:7mm) .. (br2);
  \draw[->, thick, omProp, shorten <=4pt, shorten >=3pt]
    ($(brm)+(0,-0.25)$) .. controls +(-120:8mm) and +(-60:8mm) .. (cp);}

\vspace{6mm}
\textbf{\omterm{def:g12:curly-arrows:categories}{Elimination}} (acid dehydration of ethanol, three steps):
\ce{CH3-CH2-OH -> CH2=CH2 + H2O}

\medskip
\begin{tabular}{@{}l@{}}
1.\ the oxygen takes a hydrogen \omterm{def:g9:ions:ion}{ion}: \ce{CH3-CH2-OH + H+ -> CH3-CH2-OH2+};\\
2.\ the \ce{C-O} pair leaves with water: \ce{CH3-CH2-OH2+ -> CH3-CH2+ + H2O};\\
3.\ a neighbouring \ce{C-H} pair becomes the \omterm{def:g10:lewis-and-shape:multiple-bond}{double bond}, and the
    \ce{H+} leaves:\\
\phantom{3.\ }\ce{CH3-CH2+ -> CH2=CH2 + H+}.
\end{tabular}

\medskip
{\small Three mechanisms. Orange arrows: a donor pair forms a new bond.
Blue arrows: a bond breaks, its pair going to the more electronegative
\omterm{def:g7:atoms-and-molecules:atom}{atom}. In the \omterm{def:g12:curly-arrows:categories}{elimination} the hydrogen \omterm{def:g9:ions:ion}{ion} is given back at the end.}
\end{omfigure}


\begin{remark}[Changing the skeleton or the group]\label{rem:g12:curly-arrows:skeleton}
Some reactions change only the \omterm{def:g11:functional-groups:functional-group}{functional group} and keep the carbon
skeleton: the \omterm{def:g12:curly-arrows:categories}{substitution} of bromoethane turns a \omterm{def:g11:functional-groups:halogenoalkane}{halogenoalkane} into an
\omterm{def:g11:functional-groups:alcohol}{alcohol}, two carbons before, two after. Others build or cut the
skeleton: making a new carbon--carbon bond lengthens a chain, the key
step in making a large \omterm{def:g7:atoms-and-molecules:molecule}{molecule} from small ones. Planning a \omterm{def:g10:synthesis-yield:synthesis}{synthesis}
means choosing reactions of both kinds, in the right order.
\end{remark}

\section{Exercises}

\begin{exercise}[$\star$]\label{exo:g12:curly-arrows:1}
In each species, name a donor site and, if there is one, an acceptor
site: \ce{OH-}; \ce{H+}; \ce{CH2=CH2}; \ce{CH3-Cl}; \ce{H2O}.
\end{exercise}

\begin{exercise}[$\star$]\label{exo:g12:curly-arrows:2}
\omterm{def:g12:curly-arrows:categories}{Substitution}, \omterm{def:g12:curly-arrows:categories}{addition} or \omterm{def:g12:curly-arrows:categories}{elimination}?
(a) \ce{CH3-CH2-Cl + OH- -> CH3-CH2-OH + Cl-};
(b) \ce{CH2=CH2 + H2O -> CH3-CH2-OH};
(c) \ce{CH3-CH2-OH -> CH2=CH2 + H2O};
(d) \ce{CH2=CH2 + Br2 -> CH2Br-CH2Br}.
\end{exercise}

\begin{exercise}[$\star$]\label{exo:g12:curly-arrows:3}
Where does a \omterm{def:g12:curly-arrows:curly-arrow}{curly arrow} start, and where does it end? What does a \omterm{def:g12:curly-arrows:curly-arrow}{curly
arrow} from a \omterm{def:g10:lewis-and-shape:lone-pair}{lone pair} of an oxygen \omterm{def:g7:atoms-and-molecules:atom}{atom} to a carbon \omterm{def:g7:atoms-and-molecules:atom}{atom} mean?
\end{exercise}

\begin{exercise}[$\star$]\label{exo:g12:curly-arrows:4}
In the reaction \ce{H+ + NH3 -> NH4+}, which species gives the pair?
Draw the arrow. What is the charge of the product?
\end{exercise}

\begin{exercise}[$\star$]\label{exo:g12:curly-arrows:5}
What is a \omterm{def:g12:curly-arrows:mechanism}{reaction intermediate}? Name the intermediate of the \omterm{def:g12:curly-arrows:categories}{addition} of
\ce{HBr} to ethene.
\end{exercise}

\begin{exercise}[$\star\star$]\label{exo:g12:curly-arrows:6}
Draw the \omterm{def:g12:curly-arrows:curly-arrow}{curly arrows} of the \omterm{def:g12:curly-arrows:categories}{substitution}
\ce{CH3-Br + OH- -> CH3-OH + Br-} (one step), and check the charges.
\end{exercise}

\begin{exercise}[$\star\star$]\label{exo:g12:curly-arrows:7}
In the first step of the \omterm{def:g12:curly-arrows:categories}{addition} of \ce{HBr} to ethene, why does the
\ce{H-Br} pair go to bromine rather than to hydrogen? What charge does
each product of the step carry?
\end{exercise}

\begin{exercise}[$\star\star$]\label{exo:g12:curly-arrows:8}
In the dehydration of ethanol, write the three steps and say, for each,
which species is the donor and which the acceptor. Why is the hydrogen
\omterm{def:g9:ions:ion}{ion} not written in the overall equation?
\end{exercise}

\begin{exercise}[$\star\star$]\label{exo:g12:curly-arrows:9}
The step \ce{CH3-CH2+ + H2O -> CH3-CH2-OH2+} follows the
\omterm{def:g12:curly-arrows:categories}{addition} of a hydrogen \omterm{def:g9:ions:ion}{ion} to ethene in water. Draw the arrow. What is
formed after a further loss of \ce{H+}? Write the overall equation and
give its category.
\end{exercise}

\begin{exercise}[$\star\star$]\label{exo:g12:curly-arrows:10}
Read the mechanism of the \omterm{def:g12:curly-arrows:categories}{substitution} of bromoethane: how many pairs
move? Which bond forms, which breaks? Is there an intermediate?
\end{exercise}

\begin{exercise}[$\star\star$]\label{exo:g12:curly-arrows:11}
% ledger: en:C, en:Cl, en:Br
Chloromethane reacts with the hydroxide \omterm{def:g9:ions:ion}{ion} more slowly than
bromomethane. Using electronegativities (C 2.55, Cl 3.16, Br 2.96),
explain which carbon is the more $\delta^+$. Does that explain the
speeds? (Think also of how easily the leaving \omterm{def:g9:ions:ion}{ion} takes the pair.)
\end{exercise}

\begin{exercise}[$\star\star\star$]\label{exo:g12:curly-arrows:12}
Hydrogen chloride adds to propene, \ce{CH2=CH-CH3}. In the first step
the hydrogen \omterm{def:g9:ions:ion}{ion} can bond to either carbon of the \omterm{def:g10:lewis-and-shape:multiple-bond}{double bond}, giving
two possible carbocations. Draw both, the two possible products, and
name them. (Which one forms most is studied in the Year 1 volume.)
\end{exercise}

\begin{exercise}[$\star\star\star$]\label{exo:g12:curly-arrows:13}
A student draws the first step of the \omterm{def:g12:curly-arrows:categories}{substitution} of bromoethane with an
arrow from the bromine \omterm{def:g7:atoms-and-molecules:atom}{atom} to the oxygen of \ce{OH-}. Explain the two
mistakes.
\end{exercise}

\begin{exercise}[$\star\star\star$]\label{exo:g12:curly-arrows:14}
In the formation of an \omterm{def:g11:functional-groups:carboxylic-acid}{ester} from a \omterm{def:g11:functional-groups:carboxylic-acid}{carboxylic acid} and an \omterm{def:g11:functional-groups:alcohol}{alcohol}, the
oxygen of the \omterm{def:g11:functional-groups:alcohol}{alcohol} attacks the carbon of the \ce{C=O} group. Using
\omterm{def:g11:polarity-and-cohesion:polar-bond}{partial charges}, explain why that carbon is an acceptor site and that
oxygen a donor site. Which small \omterm{def:g7:atoms-and-molecules:molecule}{molecule} is lost overall, and what
category is the whole reaction?
\end{exercise}

\begin{exercise}[$\star\star\star$]\label{exo:g12:curly-arrows:15}
Ethene reacts with bromine \ce{Br2}, a \omterm{def:g11:polarity-and-cohesion:polar-molecule}{non-polar molecule}. As a \ce{Br2}
\omterm{def:g7:atoms-and-molecules:molecule}{molecule} approaches the \omterm{def:g10:lewis-and-shape:multiple-bond}{double bond}, the \omterm{def:g9:inside-the-atom:nucleus}{electrons} of \ce{Br-Br} are
pushed away from the near bromine \omterm{def:g7:atoms-and-molecules:atom}{atom}. Explain why this creates an
acceptor site, and propose the first step of the mechanism.
\end{exercise}

\section{Problem: Making Bromoethane}

\begin{problem}[{Weekend problem --- from ten grams of ethene: how much
bromoethane?}]
\label{pb:g12:curly-arrows:1}
% ledger: aw:C, aw:H, aw:Br, en:H, en:Br, en:C
Bromoethane, a starting material for many syntheses, is made by adding
hydrogen bromide to ethene: \ce{CH2=CH2 + HBr -> CH3-CH2-Br}. A chemist
uses \qty{10.0}{g} of ethene and an excess of hydrogen bromide; the
\omterm{def:g10:synthesis-yield:yield}{yield} is \qty{75}{\percent}.

\medskip
\noindent\textbf{Part I --- The sites.}
\begin{enumerate}
  \item Draw the \omterm{def:g10:lewis-and-shape:lewis-structure}{Lewis structures} of ethene and of hydrogen bromide.
  \item Where is the donor site of ethene? Why?
  \item Compute the difference of \omterm{def:g11:polarity-and-cohesion:electronegativity}{electronegativity} of the \ce{H-Br} bond
        (H 2.20, Br 2.96). Which \omterm{def:g7:atoms-and-molecules:atom}{atom} carries $\delta^+$?
  \item Which \omterm{def:g7:atoms-and-molecules:atom}{atom} of \ce{HBr} is the acceptor site?
\end{enumerate}

\medskip
\noindent\textbf{Part II --- The mechanism.}
\begin{enumerate}[resume]
  \item Draw the first step with its two \omterm{def:g12:curly-arrows:curly-arrow}{curly arrows}.
  \item Give the two species formed and their charges.
  \item Check the conservation of charge in the first step.
  \item Draw the second step with its \omterm{def:g12:curly-arrows:curly-arrow}{curly arrow}.
  \item How many \omterm{def:g9:inside-the-atom:nucleus}{electrons} surround the positive carbon of the
        intermediate? Why is it so reactive?
\end{enumerate}

\medskip
\noindent\textbf{Part III --- What kind of reaction?}
\begin{enumerate}[resume]
  \item What is the intermediate of the mechanism?
  \item Why does it not appear in the overall equation?
  \item To which category does the reaction belong?
  \item Does the reaction change the carbon skeleton, or only the
        \omterm{def:g11:functional-groups:functional-group}{functional group}?
\end{enumerate}

\medskip
\noindent\textbf{Part IV --- How much product?}
\begin{enumerate}[resume]
  \item Compute the molar masses of ethene and of bromoethane.
  \item Compute the amount of ethene.
  \item Fill the \omterm{def:g10:reaction-progress-table:extent}{progress table}. Which \omterm{def:g7:chemical-reactions:reactant}{reactant} is limiting?
  \item Compute the maximum mass of bromoethane.
  \item Compute the mass obtained with a \omterm{def:g10:synthesis-yield:yield}{yield} of \qty{75}{\percent}.
  \item State the final answer: what mass of bromoethane does the chemist
        obtain?
\end{enumerate}
\end{problem}
